\section{Worked problems: batched steps, split counts and the lab's roofline}

\subsection{When attention outweighs the weights}
\textbf{Problem.} Llama-3-8B decodes a batch of 64 sequences, each with 4,096
tokens of context, on an H100 (3.35~TB/s), with BF16 weights and a BF16 cache.
How many bytes does one step read for the weights and for attention, and what
share is attention's? Compare batch one, and a batch of 64 with an FP8 cache.

\textbf{Solution.} The model's shape (32 layers, width 4,096, feed-forward width
14,336, vocabulary 128,256, untied embeddings) gives 8.03~billion parameters,
16.1~GB in BF16, read once per step whatever the batch: 4.8~ms at 3.35~TB/s. Its
cache is $2 \times 32 \times 8 \times 128 \times 2 = 131{,}072$~bytes per token,
so each sequence's 4,096 tokens are 0.54~GB and the batch's are 34.4~GB: 10.3~ms.
Attention is 68\% of the step's bytes, and the step takes at least 15.1~ms. At
batch one, the cache is 3\% of the bytes; the step is the weights. An FP8 cache
halves attention's share of time to 5.1~ms, 52\% of a 9.9~ms step (all derived).
The arithmetic confirms that bytes, not FLOPs, set the pace: the batch's
attention is 137~GFLOP, 0.14~ms at the H100's 989~TFLOP/s, at the 4~FLOP/byte of
equation~\eqref{eq:attn-intensity} with $g = 4$. Batching amortizes the weights,
never the cache --- which is why the chapters on cache quantization, paging and
sparse attention are about throughput, not only about memory capacity.

\subsection{How many splits a GPU needs}
\textbf{Problem.} The same model decodes one sequence at 32,768 tokens on an H100
with 132 SMs. How many independent blocks does its attention launch without
splitting, per query head and per packed key/value head? With 256-token
chunks? How much traffic do the partials add? Then answer the same questions
for a batch of 64 sequences at 4,096 tokens.

\textbf{Solution.} Unsplit, the step has 32 blocks per layer (one per query head)
or 8 (one per packed key/value head): a quarter, or a sixteenth, of the SMs,
each walking 128~MiB of keys and values alone. With 256-token chunks each head
gets 128 of them: 4,096 blocks, or 1,024 packed --- 31 or 7.8 per SM, enough to
keep every SM busy. The partials are 130 floats per (query head, chunk):
$32 \times 128 \times 520$~bytes = 2.1~MB per layer, written once and read once
by the merge, against 134~MB of cache per layer --- about 3\% more traffic
(derived). At batch 64 and 4,096 tokens the batch alone supplies
$64 \times 8 = 512$ packed blocks, 3.9 per SM, before any split. Fixed 256-token
chunks would still cut each sequence 16 ways: 8,192 blocks and 16 partials per
head to merge, 17~MB of partials per layer, written and read once beside 537~MB
of cache: 6\% more traffic for parallelism the batch already had. A fixed chunk size costs a
little at large batches in exchange for results that do not depend on the
batch. A split chosen to fill the machine for each batch saves that 6\% and
gives up the guarantee: FlashInfer's balanced scheduler, for one, picks the
largest split size that still saturates the GPU \cite{he2025nondeterminism}
(\Chref{fast-wrong-answers}).

\subsection{The lab on the roofline}
\textbf{Problem.} At 32,768 tokens the lab's step performs 268~MFLOP (16 query
heads $\times$ 32,768 tokens $\times$ 512 FLOPs) over a cache of 67~MB, counted
once. One thread runs the \texttt{packed} plan in 12.2~ms. \Chref{roofline}'s lab
measured 43.6~GB/s of reads for one thread and 60~GB/s for four. Is one thread
compute-bound or memory-bound? Predict the four-thread times of
\texttt{packed}, \texttt{packed+split} and \texttt{heads}, and compare them with
the measurements.

\textbf{Solution.} One thread computes at $268~\text{MFLOP} / 12.2~\text{ms} =
22$~GFLOP/s, while reading the cache once would take $67~\text{MB} / 43.6~\text{GB/s}
= 1.5$~ms: the thread spends eight times longer computing than it would reading,
so it is compute-bound. Four cores give about 88~GFLOP/s. The \texttt{packed} plan
has two tasks, so it can use two cores: 44~GFLOP/s, 6.1~ms predicted, 6.3
measured. \texttt{packed+split} can use all four: 3.05~ms predicted, 3.6
measured, the difference being the merges, the task queue and starting the
threads. The \texttt{heads} plan does the same arithmetic but asks for 537~MB,
each key and value eight times. If every request went to memory, it would take
$537~\text{MB} / 60~\text{GB/s} = 8.9$~ms; if none did, 3.05~ms like the packed
version. It measured 4.7--5.4~ms. In 4.7~ms four cores can draw at most 284~MB
from memory, so at least 47\% of the repeated reads came from the caches, and in
5.4~ms at least 39\% (all derived). The same arithmetic on the M1's GPU, where
llama.cpp's per-head kernel ran at close to the read-once cost, says that there
nearly all of them did.
